\usepackage{amsmath,amssymb,amsthm}

\providecommand{\lin}[1]{\ensuremath{\left\langle #1 \right\rangle}}
\providecommand{\abs}[1]{\left\lvert#1\right\rvert}
\providecommand{\norm}[1]{\left\lVert#1\right\rVert}

\providecommand{\refLE}[1]{\ensuremath{\stackrel{(\ref{#1})}{\leq}}}
\providecommand{\refEQ}[1]{\ensuremath{\stackrel{(\ref{#1})}{=}}}
\providecommand{\refGE}[1]{\ensuremath{\stackrel{(\ref{#1})}{\geq}}}
\providecommand{\refID}[1]{\ensuremath{\stackrel{(\ref{#1})}{\equiv}}}

\providecommand{\R}{\mathbb{R}} %
\providecommand{\N}{\mathbb{N}} %

\providecommand{\E}{{\mathbb E}}
\providecommand{\E}[1]{{\mathbb E}\left.#1\right. }        %
\providecommand{\Eb}[1]{{\mathbb E}\left[#1\right] }       %
\providecommand{\EE}[2]{{\mathbb E}_{#1}\left.#2\right. }  %
\providecommand{\EEb}[2]{{\mathbb E}_{#1}\left[#2\right] } %
\providecommand{\prob}[1]{{\rm Pr}\left[#1\right] }
\providecommand{\Prob}[2]{{\rm Pr}_{#1}\left[#2\right] }
\providecommand{\P}[1]{{\rm Pr}\left.#1\right. }
\providecommand{\Pb}[1]{{\rm Pr}\left[#1\right] }
\providecommand{\PP}[2]{{\rm Pr}_{#1}\left[#2\right] }
\providecommand{\PPb}[2]{{\rm Pr}_{#1}\left[#2\right] }
\providecommand{\derive}[2]{\frac{{\partial}{#1}}{\partial #2} }  %
\providecommand{\dm}{\mathrm{d}}

\DeclareMathOperator*{\argmin}{arg\,min\,}
\DeclareMathOperator*{\argmax}{arg\,max\,}
\DeclareMathOperator*{\supp}{supp}
\DeclareMathOperator*{\diag}{diag}
\DeclareMathOperator*{\Tr}{Tr}

\providecommand{\0}{\mathbf{0}}
\providecommand{\1}{\mathbf{1}}
\renewcommand{\aa}{\mathbf{a}}
\providecommand{\bb}{\mathbf{b}}
\providecommand{\cc}{\mathbf{c}}
\providecommand{\dd}{\mathbf{d}}
\providecommand{\ee}{\mathbf{e}}
\providecommand{\ff}{\mathbf{f}}
\let\ggg\gg
\renewcommand{\gg}{\mathbf{g}}
\providecommand{\gv}{\mathbf{g}}
\providecommand{\hh}{\mathbf{h}}
\providecommand{\ii}{\mathbf{i}}
\providecommand{\jj}{\mathbf{j}}
\providecommand{\kk}{\mathbf{k}}
\let\lll\ll
\renewcommand{\ll}{\mathbf{l}}
\providecommand{\mm}{\mathbf{m}}
\providecommand{\nn}{\mathbf{n}}
\providecommand{\oo}{\mathbf{o}}
\providecommand{\pp}{\mathbf{p}}
\providecommand{\qq}{\mathbf{q}}
\providecommand{\rr}{\mathbf{r}}
\renewcommand{\ss}{\mathbf{s}}
\providecommand{\tt}{\mathbf{t}}
\providecommand{\uu}{\mathbf{u}}
\providecommand{\vv}{\mathbf{v}}
\providecommand{\ww}{\mathbf{w}}
\providecommand{\xx}{\mathbf{x}}
\providecommand{\yy}{\mathbf{y}}
\providecommand{\zz}{\mathbf{z}}

\providecommand{\mA}{\mathbf{A}}
\providecommand{\mB}{\mathbf{B}}
\providecommand{\mC}{\mathbf{C}}
\providecommand{\mD}{\mathbf{D}}
\providecommand{\mE}{\mathbf{E}}
\providecommand{\mF}{\mathbf{F}}
\providecommand{\mG}{\mathbf{G}}
\providecommand{\mH}{\mathbf{H}}
\providecommand{\mI}{\mathbf{I}}
\providecommand{\mJ}{\mathbf{J}}
\providecommand{\mK}{\mathbf{K}}
\providecommand{\mL}{\mathbf{L}}
\providecommand{\mM}{\mathbf{M}}
\providecommand{\mN}{\mathbf{N}}
\providecommand{\mO}{\mathbf{O}}
\providecommand{\mP}{\mathbf{P}}
\providecommand{\mQ}{\mathbf{Q}}
\providecommand{\mR}{\mathbf{R}}
\providecommand{\mS}{\mathbf{S}}
\providecommand{\mT}{\mathbf{T}}
\providecommand{\mU}{\mathbf{U}}
\providecommand{\mV}{\mathbf{V}}
\providecommand{\mW}{\mathbf{W}}
\providecommand{\mX}{\mathbf{X}}
\providecommand{\mY}{\mathbf{Y}}
\providecommand{\mZ}{\mathbf{Z}}
\providecommand{\mLambda}{\boldsymbol{\Lambda}}
\providecommand{\mOmega}{\boldsymbol{\Omega}}
\providecommand{\mTheta}{\boldsymbol{\Theta}}
\providecommand{\mpi}{\boldsymbol{\pi}}
\providecommand{\malpha}{\boldsymbol{\alpha}}
\providecommand{\meta}{\boldsymbol{\eta}}
\providecommand{\mxi}{\boldsymbol{\xi}}
\providecommand{\mepsilon}{\boldsymbol{\epsilon}}
\providecommand{\mphi}{\boldsymbol{\phi}}
\providecommand{\momega}{\boldsymbol{\omega}}
\providecommand{\mdelta}{\boldsymbol{\delta}}
\providecommand{\mvarsigma}{\boldsymbol{\varsigma}}
\providecommand{\mnu}{\boldsymbol{\nu}}
\providecommand{\mtheta}{\boldsymbol{\theta}}
\providecommand{\mmu}{\boldsymbol{\mu}}
\providecommand{\mf}{\boldsymbol{f}}
\providecommand{\mv}{\boldsymbol{v}}
\providecommand{\mg}{\boldsymbol{g}}
\providecommand{\mmF}{\boldsymbol{F}}
\providecommand{\mmA}{\boldsymbol{A}}
\providecommand{\mmB}{\boldsymbol{B}}
\providecommand{\mmG}{\boldsymbol{G}}

\providecommand{\cA}{\mathcal{A}}
\providecommand{\cB}{\mathcal{B}}
\providecommand{\cC}{\mathcal{C}}
\providecommand{\cD}{\mathcal{D}}
\providecommand{\cE}{\mathcal{E}}
\providecommand{\cF}{\mathcal{F}}
\providecommand{\cG}{\mathcal{G}}
\providecommand{\cH}{\mathcal{H}}
\providecommand{\cI}{\mathcal{I}}
\providecommand{\cJ}{\mathcal{J}}
\providecommand{\cK}{\mathcal{K}}
\providecommand{\cL}{\mathcal{L}}
\providecommand{\cM}{\mathcal{M}}
\providecommand{\cN}{\mathcal{N}}
\providecommand{\cO}{\mathcal{O}}
\providecommand{\cP}{\mathcal{P}}
\providecommand{\cQ}{\mathcal{Q}}
\providecommand{\cR}{\mathcal{R}}
\providecommand{\cS}{\mathcal{S}}
\providecommand{\cT}{\mathcal{T}}
\providecommand{\cU}{\mathcal{U}}
\providecommand{\cV}{\mathcal{V}}
\providecommand{\cX}{\mathcal{X}}
\providecommand{\cY}{\mathcal{Y}}
\providecommand{\cW}{\mathcal{W}}
\providecommand{\cZ}{\mathcal{Z}}

\newenvironment{talign}
{\let\displaystyle\textstyle\align}
{\endalign}
\newenvironment{talign*}
{\let\displaystyle\textstyle\csname align*\endcsname}
{\endalign}



\usepackage[utf8]{inputenc}         %
\usepackage[T1]{fontenc}            %
\usepackage{url}                    %
\usepackage{booktabs}               %
\usepackage{amsfonts}               %
\usepackage{nicefrac}               %
\usepackage{microtype}              %
\usepackage{algorithm}
\usepackage{algpseudocode}  %
\algnewcommand\algorithmicinput{\textbf{Input:}}
\algnewcommand\Input{\item[\algorithmicinput]}
\renewcommand{\algorithmicensure}{\textbf{Output:}}
\usepackage{graphicx}
\usepackage{subcaption}
\usepackage[flushleft]{threeparttable}
\usepackage{float}
\usepackage{multirow}
\usepackage{xspace}
\usepackage{enumitem}
\usepackage[font=small]{caption}
\usepackage{autobreak}
\usepackage{sidecap}
\usepackage{wrapfig}
\usepackage{bbding}
\usepackage[toc, page, header]{appendix}
\usepackage{tikz}
\usepackage{pifont}
\usepackage{mdframed}
\usepackage{colortbl}

\usepackage{iclr2026_conference,times}
\usepackage[most]{tcolorbox}
\tcbuselibrary{listingsutf8}
\usepackage{mdframed}
\usepackage{hyperref}
\usepackage{url}
\usepackage{booktabs}
\usepackage{xcolor} %
\newcommand{\added}[1]{\textcolor{blue}{#1}} %
\newenvironment{addedpar}
{
  \color{blue} 
  \begingroup
}
{
  \endgroup
  \par 
}
\usepackage{graphicx} %
\usepackage{multirow}
\usepackage{booktabs}
\usepackage{array}
\usepackage{enumitem}
\usepackage{makecell}

\hypersetup{
    colorlinks=true,
    linkcolor=blue,
    citecolor=blue,
    urlcolor=blue
}

\definecolor{coral}{RGB}{255,127,80}
\definecolor{darkgreen}{RGB}{0,100,0}
\definecolor{darkyellow}{RGB}{204,153,0}
\definecolor{salmon}{RGB}{250,128,114}
\definecolor{darkred}{RGB}{150,0,0}
\newcommand{\darkredtext}[1]{{\color{darkred}#1}}

\newcommand{\secref}[1]{\hyperref[#1]{\darkredtext{Sec.~\ref*{#1}}}}
\newcommand{\thmref}[1]{\hyperref[#1]{\darkredtext{Thm.~\ref*{#1}}}}
\newcommand{\defref}[1]{\hyperref[#1]{\darkredtext{Def.~\ref*{#1}}}}
\newcommand{\propref}[1]{\hyperref[#1]{\darkredtext{Prop.~\ref*{#1}}}}
\newcommand{\assumpref}[1]{\hyperref[#1]{\darkredtext{Assump.~\ref*{#1}}}}
\newcommand{\remarkref}[1]{\hyperref[#1]{\darkredtext{Rem.~\ref*{#1}}}}
\newcommand{\hypref}[1]{\hyperref[#1]{\darkredtext{Hyp.~\ref*{#1}}}}
\newcommand{\conjref}[1]{\hyperref[#1]{\darkredtext{Conj.~\ref*{#1}}}}
\newcommand{\lemref}[1]{\hyperref[#1]{\darkredtext{Lem.~\ref*{#1}}}}
\newcommand{\corref}[1]{\hyperref[#1]{\darkredtext{Cor.~\ref*{#1}}}}
\newcommand{\noteref}[1]{\hyperref[#1]{\darkredtext{Nota.~\ref*{#1}}}}
\newcommand{\claimref}[1]{\hyperref[#1]{\darkredtext{Clm.~\ref*{#1}}}}
\newcommand{\obsref}[1]{\hyperref[#1]{\darkredtext{Obs.~\ref*{#1}}}}
\newcommand{\myalgref}[1]{\hyperref[#1]{\darkredtext{Alg.~\ref*{#1}}}}
\newcommand{\figref}[1]{\hyperref[#1]{\darkredtext{Fig.~\ref*{#1}}}}
\newcommand{\tabref}[1]{\hyperref[#1]{\darkredtext{Tab.~\ref*{#1}}}}
\newcommand{\appref}[1]{\hyperref[#1]{\darkredtext{App.~\ref*{#1}}}}

\newtheoremstyle{custom}
{1pt} %
{1pt} %
{\itshape} %
{} %
{\bfseries} %
{} %
{ } %
{\thmname{#1} \thmnumber{#2} \thmnote{(#3)} . } %

\theoremstyle{custom}

\newtheorem{innerdefinition}{Definition}
\newtheorem{innerproposition}{Proposition}
\newtheorem{innerassumption}{Assumption}
\newtheorem{innerremark}{Remark}
\newtheorem{innertheorem}{Theorem}
\newtheorem{innerhypothesis}{Hypothesis}
\newtheorem{innerconjecture}{Conjecture}
\newtheorem{innerlemma}{Lemma}
\newtheorem{innercorollary}{Corollary}
\newtheorem{innernotation}{Notation}
\newtheorem{innerclaim}{Claim}
\newtheorem{innerproblem}{Problem}
\newtheorem{innerexercise}{Exercise}
\newtheorem{innerobservation}{Observation}
\newtheorem{innernote}{Note}

\newmdenv[
    backgroundcolor=gray!10,
    linecolor=gray!100,
    linewidth=0.8pt,
    skipabove=2pt,
    skipbelow=2pt,
    innertopmargin=10pt,
    innerbottommargin=5pt,
    innerleftmargin=5pt,
    innerrightmargin=5pt,
]{definitionframe}

\newmdenv[
    backgroundcolor=blue!10,
    linecolor=blue!100,
    linewidth=0.8pt,
    skipabove=2pt,
    skipbelow=2pt,
    innertopmargin=10pt,
    innerbottommargin=5pt,
    innerleftmargin=5pt,
    innerrightmargin=5pt,
]{propositionframe}

\newmdenv[
    backgroundcolor=green!10,
    linecolor=green!100,
    linewidth=0.8pt,
    skipabove=2pt,
    skipbelow=2pt,
    innertopmargin=10pt,
    innerbottommargin=5pt,
    innerleftmargin=5pt,
    innerrightmargin=5pt,
]{assumptionframe}

\newmdenv[
    backgroundcolor=yellow!10,
    linecolor=yellow!100,
    linewidth=0.8pt,
    skipabove=2pt,
    skipbelow=2pt,
    innertopmargin=10pt,
    innerbottommargin=5pt,
    innerleftmargin=5pt,
    innerrightmargin=5pt,
]{remarkframe}

\newmdenv[
    backgroundcolor=red!10,
    linecolor=red!100,
    linewidth=0.8pt,
    skipabove=2pt,
    skipbelow=2pt,
    innertopmargin=10pt,
    innerbottommargin=5pt,
    innerleftmargin=5pt,
    innerrightmargin=5pt,
]{theoremframe}

\newmdenv[
    backgroundcolor=purple!10,
    linecolor=purple!100,
    linewidth=0.8pt,
    skipabove=2pt,
    skipbelow=2pt,
    innertopmargin=10pt,
    innerbottommargin=5pt,
    innerleftmargin=5pt,
    innerrightmargin=5pt,
]{hypothesisframe}

\newmdenv[
    backgroundcolor=orange!10,
    linecolor=orange!100,
    linewidth=0.8pt,
    skipabove=2pt,
    skipbelow=2pt,
    innertopmargin=10pt,
    innerbottommargin=5pt,
    innerleftmargin=5pt,
    innerrightmargin=5pt,
]{conjectureframe}

\newmdenv[
    backgroundcolor=cyan!10,
    linecolor=cyan!100,
    linewidth=0.8pt,
    skipabove=2pt,
    skipbelow=2pt,
    innertopmargin=10pt,
    innerbottommargin=5pt,
    innerleftmargin=5pt,
    innerrightmargin=5pt,
]{lemmaframe}

\newmdenv[
    backgroundcolor=magenta!10,
    linecolor=magenta!100,
    linewidth=0.8pt,
    skipabove=2pt,
    skipbelow=2pt,
    innertopmargin=10pt,
    innerbottommargin=5pt,
    innerleftmargin=5pt,
    innerrightmargin=5pt,
]{corollaryframe}

\newmdenv[
    backgroundcolor=pink!10,
    linecolor=pink!100,
    linewidth=0.8pt,
    skipabove=2pt,
    skipbelow=2pt,
    innertopmargin=10pt,
    innerbottommargin=5pt,
    innerleftmargin=5pt,
    innerrightmargin=5pt,
]{notationframe}

\newmdenv[
    backgroundcolor=violet!10,
    linecolor=violet!100,
    linewidth=0.8pt,
    skipabove=2pt,
    skipbelow=2pt,
    innertopmargin=10pt,
    innerbottommargin=5pt,
    innerleftmargin=5pt,
    innerrightmargin=5pt,
]{claimframe}

\newmdenv[
    backgroundcolor=salmon!10,
    linecolor=salmon!100,
    linewidth=0.8pt,
    skipabove=2pt,
    skipbelow=2pt,
    innertopmargin=10pt,
    innerbottommargin=5pt,
    innerleftmargin=5pt,
    innerrightmargin=5pt,
]{problemframe}

\newmdenv[
    backgroundcolor=lavender!10,
    linecolor=lavender!100,
    linewidth=0.8pt,
    skipabove=2pt,
    skipbelow=2pt,
    innertopmargin=10pt,
    innerbottommargin=5pt,
    innerleftmargin=5pt,
    innerrightmargin=5pt,
]{observationframe}


\newenvironment{definition}
{\begin{definitionframe}\begin{innerdefinition}}
            {\end{innerdefinition}\end{definitionframe}}

\newenvironment{proposition}
{\begin{propositionframe}\begin{innerproposition}}
            {\end{innerproposition}\end{propositionframe}}

\newenvironment{assumption}
{\begin{assumptionframe}\begin{innerassumption}}
            {\end{innerassumption}\end{assumptionframe}}

\newenvironment{remark}
{\begin{remarkframe}\begin{innerremark}}
            {\end{innerremark}\end{remarkframe}}

\newenvironment{theorem}
{\begin{theoremframe}\begin{innertheorem}}
            {\end{innertheorem}\end{theoremframe}}

\newenvironment{hypothesis}
{\begin{hypothesisframe}\begin{innerhypothesis}}
            {\end{innerhypothesis}\end{hypothesisframe}}

\newenvironment{conjecture}
{\begin{conjectureframe}\begin{innerconjecture}}
            {\end{innerconjecture}\end{conjectureframe}}

\newenvironment{lemma}
{\begin{lemmaframe}\begin{innerlemma}}
            {\end{innerlemma}\end{lemmaframe}}

\newenvironment{corollary}
{\begin{corollaryframe}\begin{innercorollary}}
            {\end{innercorollary}\end{corollaryframe}}

\newenvironment{notation}
{\begin{notationframe}\begin{innernotation}}
            {\end{innernotation}\end{notationframe}}

\newenvironment{claim}
{\begin{claimframe}\begin{innerclaim}}
            {\end{innerclaim}\end{claimframe}}

\newenvironment{problem}
{\begin{problemframe}\begin{innerproblem}}
            {\end{innerproblem}\end{problemframe}}

\newenvironment{observation}
{\begin{observationframe}\begin{innerobservation}}
            {\end{innerobservation}\end{observationframe}}


\newcommand{\TBD}{\textcolor{red}{TBD}}
\newcommand{\method}{\textsc{SupervisorAgent}\xspace} %
\newcommand{\methods}{\method-\{T,\,S\}\xspace} %
\newcommand{\tip}[1]{{\color{red}{[#1]}}}
\newcommand{\sg}{\mathop{{\color{red}\operatorname{sg}}}\nolimits}
\newcommand{\dg}[1]{\textcolor{darkgreen}{$^{\downarrow#1}$}}
\newcommand{\up}[1]{\textcolor{red}{$^{\uparrow#1}$}}
\newcommand{\rp}{$\boldsymbol{\color{red!70!black}\oplus}$}
\newcommand{\graycell}[1]{\cellcolor[gray]{0.9}#1}
\newcommand{\IPC}{\texttt{IPC}\xspace}
\newcommand{\SSL}{\texttt{SSL}\xspace}
\newcommand{\CMB}{\textcolor{green}{\CheckmarkBold}}
\newcommand{\Tool}{(\raisebox{-0.5ex}{\includegraphics[height=1em]{sources/figures/btool.png}})\xspace}
\newcommand{\Runer}{(\raisebox{-0.5ex}{\includegraphics[height=1em]{sources/figures/flash.png}})\xspace}
\newcommand{\XXX}{\tip{XXX~\citep{template}}}


\usepackage[textwidth=3.5cm]{todonotes}
\newcommand{\tao}[1]{\todo[color=blue!20,size=\footnotesize,caption={}]{T: #1}{}}
\newcommand{\peng}[1]{\todo[color=green!20,size=\footnotesize,caption={}]{S: #1}{}}
\newcommand{\shang}[1]{\todo[color=orange!20,size=\footnotesize,caption={}]{Y: #1}{}}

\newcommand{\itao}[1]{\todo[inline,color=blue!20,size=\footnotesize,caption={}]{T: #1}{}}
\newcommand{\ipeng}[1]{\todo[inline,color=green!20,size=\footnotesize,caption={}]{S: #1}{}}
\newcommand{\ishang}[1]{\todo[inline,color=orange!20,size=\footnotesize,caption={}]{Y: #1}{}}

\usepackage{newunicodechar}
\newunicodechar{↑}{$\uparrow$}
\newunicodechar{↓}{$\downarrow$}


\newcommand*\cw[1]{\tikz[baseline=(char.base)]{
        \node[shape=circle,fill=white,draw,inner sep=0.5pt,scale=0.8,line width=0.8pt] (char) {\textcolor{black}{#1}};}}

\newcommand*\cb[1]{\tikz[baseline=(char.base)]{
        \node[shape=circle,fill=black,draw,inner sep=0.5pt,scale=0.8,line width=0.8pt] (char) {\textcolor{white}{#1}};}}
